\section{Conclusiones}
Esta revisión tuvo como objetivo sintetizar los enfoques metodológicos empleados, resultados educativos reportados e identificar los principales desafíos y oportunidades. Los hallazgos evidencian que la inteligencia artificial aplicada a lenguas de bajos recursos se encuentra en crecimiento, aunque todavía presenta una producción científica limitada y dispersa. No obstante, se identificó un incremento en las publicaciones del año 2024 al 2025. Asimismo, la mayor parte de las publicaciones corresponde a investigaciones experimentales y propuestas de modelos las cuales representan el 66.7\% del total de la muestra. 

En cuanto a aplicaciones educativas, los resultados muestran que aplicaciones basadas en NLP contribuyen positivamente a integrar la lengua materna en entornos de aprendizaje con resultados más satisfactorios que los de una enseñanza tradicional. Por otro lado, entre las estrategias más utilizadas para enfrentar la falta de corpus y mejorar el desempeño de los modelos, se identifica el uso de descripciones lingüísticas, diccionarios mediante \textit{in-context learning} y técnicas de meta-aprendizaje. A su vez, otro hallazgo relevante es la existencia de brechas significativas en cuanto al rendimiento de herramientas basadas en IA entre lenguas de alto recurso y lenguas de bajos recursos, se reportaron distintas limitantes en tareas de comprensión semántica, traducción automática, seguimiento de instrucciones, seguridad y confiabilidad de los modelos.
En relación con las estrategias metodológicas, se evidenció que los enfoques más destacables combinan técnicas algorítmicas con metodologías participativas, como el \textit{transfer learning}, modelos multilingües, aumento de datos, meta-aprendizaje, arquitecturas híbridas y la inyección de conocimiento lingüístico. Sin embargo, también se resalta la importancia de involucrar a la población de estudio en el diseño y evaluación de las herramientas a fin de evitar la distorsión del conocimiento, respetar la soberanía sobre los datos y asegurar que las soluciones respondan a sus necesidades.
Finalmente, esta revisión sistemática de la literatura identifica desafíos que deben ser abordados en futuras investigaciones, entre ellas se encuentran la escasez de corpus digitalizados, dependencia excesiva de métricas como BLEU, TER o WER, la falta de inclusión de más de una variante dialectal, ausencia de ortografías estandarizadas y la complejidad morfofonémica. Aunque la inteligencia artificial representa una vía prometedora para mejorar la accesibilidad educativa en lenguas de bajos recursos, su implementación requiere enfoques más participativos que involucren a la población de estudio, prioricen sus necesidades y también tener en cuenta la variedad dialectal. La participación de hablantes nativos podría contribuir a crear mejores datos de entrenamiento y evitar errores culturales o de significado.



% Las conclusiones deben obtenerse, por tanto, a partir de algo más que de los simples datos registrados. De hecho, unos datos o resultados pueden tener un sentido u otro y, por tanto, pueden llevar a unas conclusiones y otras, dependiendo del marco conceptual que justifica la investigación, de la metodología seguida, de los objetivos propuestos, etc.
%========================================================
